% Figure: the Joule-Thomson inversion curve in reduced coordinates. Inside
% the dome-shaped inversion curve the JT coefficient is positive and a
% throttle cools the gas (the liquefaction region); outside it the throttle
% warms the gas. Schematic, of the van der Waals inversion shape, with the
% low-pressure maximum inversion temperature marked at T_r = 27/4 approx 6.75
% clamped to the plotted range, and the two working reduced isenthalps drawn.
\begin{figure}[htbp]
\centering
\begin{tikzpicture}
\begin{axis}[
  width=12cm, height=8cm,
  xlabel={reduced pressure $p_r=p/p_c$},
  ylabel={reduced temperature $T_r=T/T_c$},
  xmin=0, xmax=12, ymin=0, ymax=7.5,
  axis lines=left,
  tick label style={font=\scriptsize},
  label style={font=\small},
  legend pos=north east,
  legend style={font=\scriptsize},
  clip=true,
]
% Inversion curve, schematic dome. Low-pressure intercept near T_r ~ 6.75,
% peak folded over, low-T_r intercept near 0.8. Drawn as a smooth closed arc.
\addplot[engblue, very thick, smooth] coordinates {
  (0,6.7) (1.5,6.2) (3.0,5.4) (4.5,4.4) (6.0,3.4) (7.0,2.6)
  (7.4,2.0) (7.0,1.5) (6.0,1.1) (4.5,0.9) (3.0,0.82) (1.5,0.80) (0,0.80)
};
\addlegendentry{inversion curve $\mu_{JT}=0$}
% cooling region label (inside)
\node[font=\scriptsize\itshape, engblue] at (axis cs:2.6,3.2)
  {$\mu_{JT}>0$: throttle cools};
% warming region label (outside, upper right)
\node[font=\scriptsize\itshape, enggray] at (axis cs:9.0,5.5)
  {$\mu_{JT}<0$: throttle warms};
% a representative isenthalp passing through the dome (cooling on the way in)
\addplot[engorange, thick, dashed, smooth] coordinates {
  (11,5.9) (9,5.2) (7,4.3) (5,3.4) (3,2.7) (1,2.2) (0,2.05)
};
\addlegendentry{isenthalp $h=\text{const}$}
% maximum inversion temperature marker on the T_r axis
\addplot[only marks, mark=*, engred, mark size=2pt] coordinates {(0,6.7)};
\node[anchor=west, font=\scriptsize, engred] at (axis cs:0.2,6.9)
  {$T_{\text{inv}}^{\max}\!\approx\!\tfrac{27}{4}T_c$ (low-$p$ limit)};
\end{axis}
\end{tikzpicture}
\caption[The Joule-Thomson inversion curve]{The Joule-Thomson inversion
curve in reduced coordinates, the locus where $\mu_{JT}=(\partial T/\partial
p)_h=0$. Inside the dome the coefficient is positive and a pressure drop
cools the gas, the region a Linde liquefier throttles within; outside it the
coefficient is negative and a pressure drop warms the gas, which is why a
high-pressure hydrogen leak, whose inversion temperature sits below room
temperature, self-heats at a pinhole while a throttled carbon-dioxide leak
frosts. An isenthalp (dashed) crossing into the dome cools as it drops in
pressure. The maximum inversion temperature in the van der Waals low-pressure
limit is $T_{\text{inv}}^{\max}=2a/(Rb)=(27/4)\,T_c$, placing nitrogen's near
$850\,\si{\kelvin}$ and hydrogen's near $220\,\si{\kelvin}$. The curve is
schematic, of the van der Waals shape.}
\label{fig:vol04ch02:joule-thomson-inversion}
\end{figure}
